% GENERATED FILE. Edit canonical lessons, then run npm run generate:books.
% canonical-source-hash: fnv1a64:4f796dc3246bb507
% canonical-lessons: TA-S130-letter-nga, TA-S131-letter-lla, TA-S132-vowel-sign-ii, TA-S133-vowel-sign-e

\chapter{Letters From You and Your Name}
\label{ch:ta-letters-from-you-and-your-name}
\hlchaptermodality{110}

\begin{chapteropening}
\textbf{By the end of this chapter:} \emph{I can write the letters nga and lla and the vowel signs ii and e, and find each one in niingal or peyar.}

Complete the last lesson of chapter 110: i can write the letters nga and lla and the vowel signs ii and e, and find each one in niingal or peyar.
\end{chapteropening}

\section[ṅa]{\ta{ங} — the nasal before the \ta{க} in \ta{நீங்கள்}}

\label{lesson:TA-S130-letter-nga}



\begin{quote}
Before the new one: read \ta{◌ை} and \ta{ஸ} aloud, and write each once.

Now say \textbf{\ta{நீங்கள்}} (\emph{nīṅgaḷ}), \textquotedblleft{}you (respectful)\textquotedblright{}. You have read it for a long time. This lesson takes one character out of it and puts it in your hand.
\end{quote}

\subsection*{Script you'll notice: \ta{ங}}
\textbf{\ta{ங}} — \emph{ṅa}.

It is the nasal in the middle of \textbf{\ta{நீங்கள்}} (\emph{nīṅgaḷ}), the respectful \textquotedblleft{}you\textquotedblright{}. Say \emph{nīṅgaḷ} slowly: the \emph{ṅ} is made at the back of the mouth, where \emph{k} is made, which is why \ta{ங} nearly always stands right before \ta{க}.

\subsection*{Writing: \ta{ங}}
\hlblockfigure{\detokenize{figures/TA-S130-letter-nga-filmstrip.pdf}}{How \ta{ங} is written, stroke by stroke}

\begin{itemize}
\raggedright
  \item \textbf{1.} start here: draw the detached upright down, then lift
  \item \textbf{2.} put the pen down again and climb the tall left body
  \item \textbf{3.} without lifting, carry the top bar right and return inward
  \item \textbf{4.} without lifting, descend into the rounded inner turn
  \item \textbf{5.} without lifting, carry the low bar to the right
  \item \textbf{6.} without lifting, return left and finish up the inner stem — and only now lift
\end{itemize}

\textbf{Pen lifts: 1.}

\begin{quote}
Stroke order is one attested teaching order, not a national standard — Tamil handwriting is taught with school-to-school variation. Source: Sankaran Radhakrishnan, Tamil Script Learners Manual, Appendix I: Hand-movements, Frame 2, \ta{ங} (University of Texas at Austin), p. 191.
\end{quote}

\subsection*{Your turn}
\begin{itemize}
\raggedright
  \item \emph{Look:} at \textbf{\ta{நீங்கள்}} and point at \ta{ங} inside it
  \item \emph{Trace it:} \ta{ங} three times, saying \emph{ṅa} as you finish each one
  \item \emph{Write it:} \ta{ஸ}, then \ta{ங}, from memory
  \item \emph{Read it:} \textbf{\ta{நீங்கள்}} aloud once more
\end{itemize}

\begin{tcolorbox}[breakable,colback=teal!4,colframe=teal!35!black,title={Before you move on}]
Which character is this — \ta{ங}? (\textbf{ṅa}.) Which word did it come from? (\textbf{\ta{நீங்கள்}}, \emph{nīṅgaḷ}, \textquotedblleft{}you (respectful)\textquotedblright{}.)
\end{tcolorbox}
\section[ḷa]{\ta{ள} — the l at the end of \ta{நீங்கள்}}

\label{lesson:TA-S131-letter-lla}



\begin{quote}
Before the new one: read \ta{ஸ} and \ta{ங} aloud, and write each once.

Now say \textbf{\ta{நீங்கள்}} (\emph{nīṅgaḷ}), \textquotedblleft{}you (respectful)\textquotedblright{}. You have read it for a long time. This lesson takes one character out of it and puts it in your hand.
\end{quote}

\subsection*{Script you'll notice: \ta{ள}}
\textbf{\ta{ள}} — \emph{ḷa}.

It closes \textbf{\ta{நீங்கள்}} (\emph{nīṅgaḷ}). This is the Tamil \emph{l} made with the tongue curled back. You already write \textbf{\ta{ல}} (\emph{la}) and \textbf{\ta{ழ}} (\emph{ḻa}); \ta{ள} is the third \emph{l}, and the dot under the romanization (\emph{ḷ}) is how the book tells them apart.

\subsection*{Writing: \ta{ள}}
\hlblockfigure{\detokenize{figures/TA-S131-letter-lla-filmstrip.pdf}}{How \ta{ள} is written, stroke by stroke}

\begin{itemize}
\raggedright
  \item \textbf{1.} start here: curl around the large left loop
  \item \textbf{2.} without lifting, turn inward through the loop
  \item \textbf{3.} without lifting, descend around the inner curl, then lift
  \item \textbf{4.} put the pen down again and rise on the adjoining upright
  \item \textbf{5.} without lifting, carry the top bar right, then lift
  \item \textbf{6.} put the pen down again and descend the separate right upright — and only now lift
\end{itemize}

\textbf{Pen lifts: 2.}

\begin{quote}
Stroke order is one attested teaching order, not a national standard — Tamil handwriting is taught with school-to-school variation. Source: Sankaran Radhakrishnan, Tamil Script Learners Manual, Appendix I: Hand-movements, Frame 12, \ta{ள} (University of Texas at Austin), p. 195.
\end{quote}

\subsection*{Your turn}
\begin{itemize}
\raggedright
  \item \emph{Look:} at \textbf{\ta{நீங்கள்}} and point at \ta{ள} inside it
  \item \emph{Trace it:} \ta{ள} three times, saying \emph{ḷa} as you finish each one
  \item \emph{Write it:} \ta{ங}, then \ta{ள}, from memory
  \item \emph{Read it:} \textbf{\ta{நீங்கள்}} aloud once more
\end{itemize}

\begin{tcolorbox}[breakable,colback=teal!4,colframe=teal!35!black,title={Before you move on}]
Which character is this — \ta{ள}? (\textbf{ḷa}.) Which word did it come from? (\textbf{\ta{நீங்கள்}}, \emph{nīṅgaḷ}, \textquotedblleft{}you (respectful)\textquotedblright{}.)
\end{tcolorbox}
\section[ī]{\ta{◌ீ} — the long i on \ta{நீ}}

\label{lesson:TA-S132-vowel-sign-ii}



\begin{quote}
Before the new one: read \ta{ங} and \ta{ள} aloud, and write each once.

Now say \textbf{\ta{நீ}} (\emph{nī}), \textquotedblleft{}you (familiar)\textquotedblright{}. You have read it for a long time. This lesson takes one character out of it and puts it in your hand.
\end{quote}

\subsection*{Script you'll notice: \ta{◌ீ}}
\textbf{\ta{◌ீ}} — \emph{ī}.

It is the sign on \textbf{\ta{நீ}} (\emph{nī}), \textquotedblleft{}you\textquotedblright{}. You already write its short partner \ta{◌ி} (\emph{i}), the hook from \ta{நன்றி}. The long sign \ta{◌ீ} ends in a curl instead of an open hook. One curl more, one beat longer.

\subsection*{Writing: \ta{◌ீ}}
Put your pen on \ta{◌ீ} and follow its line. Copy the shape you can see — slowly, and larger than it is printed.

\begin{quote}
This book does not yet tell you \textbf{where to start the character or which way to travel}. That is a real thing, taught with real variation from school to school, and it is not written down here until it can be written down with a source. Copying what is in front of you needs no such source, and it is how the shape gets into your hand in the meantime.
\end{quote}

\subsection*{Your turn}
\begin{itemize}
\raggedright
  \item \emph{Look:} at \textbf{\ta{நீ}} and point at \ta{◌ீ} inside it
  \item \emph{Trace it:} \ta{◌ீ} three times, saying \emph{ī} as you finish each one
  \item \emph{Write it:} \ta{ள}, then \ta{◌ீ}, from memory
  \item \emph{Read it:} \textbf{\ta{நீ}} aloud once more
\end{itemize}

\begin{tcolorbox}[breakable,colback=teal!4,colframe=teal!35!black,title={Before you move on}]
Which character is this — \ta{◌ீ}? (\textbf{ī}.) Which word did it come from? (\textbf{\ta{நீ}}, \emph{nī}, \textquotedblleft{}you (familiar)\textquotedblright{}.)
\end{tcolorbox}
\section[e]{\ta{◌ெ} — the short e in \ta{பெயர்}}

\label{lesson:TA-S133-vowel-sign-e}



\begin{quote}
Before the new one: read \ta{ள} and \ta{◌ீ} aloud, and write each once.

Now say \textbf{\ta{பெயர்}} (\emph{peyar}), \textquotedblleft{}name\textquotedblright{}. You have read it for a long time. This lesson takes one character out of it and puts it in your hand.
\end{quote}

\subsection*{Script you'll notice: \ta{◌ெ}}
\textbf{\ta{◌ெ}} — \emph{e}.

It is the sign on the first syllable of \textbf{\ta{பெயர்}} (\emph{peyar}), \textquotedblleft{}name\textquotedblright{}. Like \ta{◌ை}, it is written on the \textbf{left} of its consonant and sounded after it. You already know its long partner \ta{◌ே} (\emph{ē}) from \ta{பேசு}. Put the two side by side and compare the shapes: that difference is the difference between \emph{pe} and \emph{pē}.

\subsection*{Writing: \ta{◌ெ}}
Put your pen on \ta{◌ெ} and follow its line. Copy the shape you can see — slowly, and larger than it is printed.

\begin{quote}
This book does not yet tell you \textbf{where to start the character or which way to travel}. That is a real thing, taught with real variation from school to school, and it is not written down here until it can be written down with a source. Copying what is in front of you needs no such source, and it is how the shape gets into your hand in the meantime.
\end{quote}

\subsection*{Your turn}
\begin{itemize}
\raggedright
  \item \emph{Look:} at \textbf{\ta{பெயர்}} and point at \ta{◌ெ} inside it
  \item \emph{Trace it:} \ta{◌ெ} three times, saying \emph{e} as you finish each one
  \item \emph{Write it:} \ta{◌ீ}, then \ta{◌ெ}, from memory
  \item \emph{Read it:} \textbf{\ta{பெயர்}} aloud once more
\end{itemize}

\begin{tcolorbox}[breakable,colback=teal!4,colframe=teal!35!black,title={Before you move on}]
Which character is this — \ta{◌ெ}? (\textbf{e}.) Which word did it come from? (\textbf{\ta{பெயர்}}, \emph{peyar}, \textquotedblleft{}name\textquotedblright{}.)
\end{tcolorbox}
